%%%%%%%%%%%%%%%%%%%%%%%%%%%%%%%%%%%%%%%%%
% University of Amsterdam master thesis title page 
% LaTeX Template
%
% Version 1.2 (14/02/14) (fixed cursive mode under titlepage)
%
% This template was made for the UvA by Ludo Nieuwenhuizen
% 
% 
% Instructions for using this template:
%
% I've done my best to make this template self-explanatory. The only thing (apart from
% finishing your thesis) is that wherever you see the percentage sign after some LaTeX
% commands, you have to insert some text (except for this
% introductory part, of course)l this is also explicitly stated.
%
% This title page can be compiled as is. This is not useful for 
% including it in another document. To do this, you have two options: 
%
% 1) Copy/paste everything between \begin{document} and \end{document} 
% starting at \begin{titlepage} and paste this into another LaTeX file where you 
% want your title page.
% OR
% 2) Remove everything outside the \begin{titlepage} and \end{titlepage} and 
% move this file to the same directory as the LaTeX file you wish to add it to. 
% Then add \input{./titlepage_master_thesis.tex} to your LaTeX file where you want your
% title page.
%
% The layout is quite vulnerable for changes. If you for instance remove the logo at the
% bottom, the last line 'Department Research institute...' will get higher up the page. In this 
% case it is safer to insert a blank image or fix it in a different way.
%
% Any questions can be sent to me (L.G.Nieuwenhuizen@gmail.com)
%%%%%%%%%%%%%%%%%%%%%%%%%%%%%%%%%%%%%%%%%

\documentclass[11pt,twoside,a4paper,fleqn]{report}

\usepackage{graphicx}
\usepackage{graphics}
\usepackage[english]{babel}
\usepackage[hcentering,bindingoffset=8mm]{geometry}
\usepackage[utf8]{inputenc}
% \usepackage[en-US]{datetime2}
\usepackage{amssymb}
\usepackage{amsmath}
\usepackage{natbib}
\usepackage[colorlinks=true, allcolors=blue]{hyperref}
\usepackage{caption}
\usepackage{tocloft}
\usepackage{amsfonts}
\usepackage[section]{placeins}
\usepackage{fancyhdr}
% \usepackage{siunitx}
\usepackage{subcaption}
\usepackage{listings}
\usepackage{acro}
% \usepackage{abbriv}
\usepackage{longtable}
\usepackage{array}
\usepackage{tablefootnote}
\usepackage{epigraph}
\usepackage[capitalise]{cleveref}
\usepackage{comment}
\usepackage{float}

\graphicspath{{folder/}{../figures/}}

\pagestyle{fancy}
\fancyhead[RO]{\fancyplain{}{\nouppercase{\leftmark}}}
\renewcommand\sectionmark[1]{\markboth{\MakeUppercase{#1}}{}}
\fancyhead[LO]{}
\fancyhead[LE]{\fancyplain{}{\nouppercase{\leftmark}}}
\renewcommand\sectionmark[1]{\markboth{\MakeUppercase{#1}}{}}
\fancyhead[RE]{}
\lfoot{}
\cfoot{\fancyplain{}{\thepage}}
\rfoot{}

% \sisetup{separate-uncertainty, multi-part-units = brackets}
% \DeclareSIUnit\parsec{pc}
% \DeclareSIUnit\photons{ph}
% \def\kms{km~s$^{-1}$}

\newcolumntype{L}[1]{>{\raggedright\let\newline\\\arraybackslash\hspace{0pt}}m{#1}}

%\setlength{\headheight}{15pt}

%\parskip = \baselineskip

\captionsetup[table]{name=Table,labelfont={sc,footnotesize},textfont=footnotesize,labelsep=none}
\captionsetup[figure]{name=Fig.,labelfont={sc,small},textfont=footnotesize,labelsep=endash}
\numberwithin{equation}{chapter}
\numberwithin{figure}{chapter}

\newcommand{\vdag}{(v)^\dagger}
\newcommand\aastex{AAS\TeX}
\newcommand\latex{La\TeX}

\setlength{\parindent}{0 cm}
\pagenumbering{gobble}

\acsetup{list-style=longtable}

\begin{document}
\begin{titlepage}

\newcommand{\HRule}{\rule{\linewidth}{0.8mm}}
\center
 \vspace*{0.5cm}  % Play around with this as you want
% ------
% Heading20
% ------
\raisebox{0.05cm}[0pt][0pt]{\includegraphics[width=2.0cm]{Thesis/logos/UvA_logo.png}}
\raisebox{0.7cm}[0pt][0pt]{\textsc{\Huge University of Amsterdam}}
\raisebox{-1.85cm}[0pt][0pt]{\includegraphics[width=7.0cm]{Thesis/logos/VUlogo.png}}\\[2.0 cm]

\Large{\textbf{MSc Physics and Astronomy}}\\% Study discipline
\Large{Track: Astronomy \& Astrophysics}\\[0.7cm] % Master track name
\textsc{\Large \textbf{Master Thesis}}\\[0.2cm]

% -----
% Title
% -----

\HRule \\[0.3cm]

{ \huge \bfseries The evolution of TIC 470710327 hierarchical triple system with a Roche lobe filling outer star }\\[0.8cm] % title of thesis
%{\Large \bfseries Subtitle of Thesis\\Can use two lines} % subtitle of thesis

\HRule \\[0.7cm]
 
% -----
% Details
% -----

{\Large \emph by}\\[0.6cm]
{\Large \bfseries Adam Parkosidis\\ % student name
13950142}\\[0.4cm] %student ID
%\DTMlangsetup{showdayofmonth=false}
{\large  \emph{\today}}\\ % month + year in which the thesis was concluded
%\DTMlangsetup{showdayofmonth=true}
{\large  \emph{60 ECTS}}\\ % 'this many' ects that are rewarded for the thesis
{\large  \emph{Start date - End date}}\\[1.8cm] % period in which the research was carried out

% -----
% Supervisor, etc.
% -----

\begin{minipage}{0.4\textwidth}
\begin{flushleft} \large
{\large \emph{Supervisors:}}\\

\large{Dr Silvia Toonen}\\ % Supervisor's Name
\large{Dr Philipp Moesta}
\end{flushleft}
\end{minipage}
~
\begin{minipage}{0.4\textwidth}
\begin{flushright} \large
\emph{Examiners} \\
\large{Dr. Silvia Toonen} \\
\large{Dr. Oliver Porth} % Examiner's Name
\end{flushright}
\end{minipage}\\


% -----
% Logo and Department
% -----

\raisebox{-138pt}[0pt][0pt]{\includegraphics[width=5.5cm]{Thesis/logos/api_logo.pdf}}\\ % Call your institute logo "institute.jpg", or be creative.

% \raisebox{-138pt}[0pt][0pt]{\large{Anton Pannekoek Institute for Astronomy}} %name of the department or institute or company


% -----
% That was easy, right?
% -----

\vfill 

\end{titlepage}

\newpage
\chapter*{Abstract}

Mass transfer in hierarchical triple systems, and more specifically, \ac{rlof} by the outer star, is expected to be considerably different from mass transfer in an ordinary binary system. The mass cannot be simply accreted by the inner objects, but may instead form a circumbinary disk or be expelled via a slingshot effect due to the inner binary rotation. Stellar evolution, gravitational dynamics, and hydrodynamics all play important roles in the process. In the first part, we create 3D hydrodynamical models of post main-sequence stars based on detailed 1D stellar evolution models. In the second part, we use AMUSE to couple hydrodynamics with high accuracy gravitational integrators and solve these physical processes in a self-consistent  manner. Hence, we simulate the phase of mass transfer in a hierarchical triple system in which the tertiary star will overfill its Roche lobe before any of the inner stars leave the main sequence. We encounter a fairly non-conservative mass transfer, and while we quantify its impact on the inner and outer orbits, predicting the end of the mass transfer phase and the appearance of the resulting system is difficult. However, we provide some preliminary estimations of the system's accretion efficiency and the amount of angular momentum lost. Finally,
we speculate that the formation of a circumbinary disk around the inner binary probably leads to significantly more conservative mass transfer. 





\mbox{}


\newpage
\pagenumbering{roman}
\setcounter{page}{1}
{
  \hypersetup{linkcolor=black}
\tableofcontents

\begin{comment}
\newpage
\thispagestyle{empty}
\mbox{}
\newpage



\listoffigures

\newpage
\thispagestyle{empty}
\mbox{}
\newpage

\listoftables

\newpage
\thispagestyle{empty}
\mbox{}
\newpage

\mbox{}
\thispagestyle{empty}
\printacronyms[include-classes=abbrev,name=List of Abbreviations]

\end{comment}

\newpage
\thispagestyle{empty}
\mbox{}
\newpage

\pagenumbering{arabic}
\setcounter{page}{1}

}

\parskip = 2mm

\chapter{Introduction}

\epigraph{The universe must be full of voices, calling from star to star in a myriad tongues. One day we shall join that cosmic conversation}{Arthur C. Clarke}

Stellar formation and evolution is one of the oldest and most studied fields in astrophysics. In their attempt to achieve hydrostatic and thermal equilibrium, stars generate temperatures and pressures that allow for nuclear burning, and thus the emission of the starlight that we see. The cycles of nuclear burning and fuel exhaustion regulate the evolution of a star and set the various phases during the stellar lifetime.

The evolution of a star in isolation, namely single star evolution, is predominantly determined by the stellar mass. As a result, stars with $M \leq 2$ M$_{\odot}$, $2 < M \leq 8$ M$_{\odot}$ and $M > 8$ M$_{\odot}$ are classified as low-, intermediate-, and high-mass, respectively. Furthermore, the evolution is only slightly affected by the initial chemical composition.  On the other hand, a number of internal mixing process can significantly alter the evolution, particularly of massive ($M>8$ M$_{\odot}$), stars \citep{langer2012presupernova}. Core overshooting, semiconvection, and rotationally induced mixing are the most important internal mixing process \citep{schootemeijer2019constraining} and are still poorly understood. 


\section{Evolution of Massive stars}

\section{Hierarchical Triple Star Systems}



Although the fundamentals of single star and binary evolution have long been acknowledged \citep{postnov2014evolution,toonen2014popcorn}, the long-term evolution of stellar triples remains unknown. In the simple case, stable triple systems, which are hierarchical, namely, consist of an inner and an outer binary orbit, i.e., the tertiary (third object). The secular evolution of such systems is the modification of orbital elements over timescales substantially larger than the system's dynamical timescale. Hence, the presence of the outer star has no influence on the history of the inner binary and the evolution of the inner binary and the tertiary can be discussed independently. In other cases, a third star in orbit around a binary system can drastically influence the development the system via interactions, which influence the orbital elements of the inner and outer orbit through changes in energy and angular momentum. Consequently, triple stellar evolution processes which are unique to systems with multiplicities of higher orders than binaries can alter stable systems to become unstable.



\newpage

\bibliographystyle{plainnat}
\bibliography{refs}

%## IMPORTANT#############################################
\acuseall
%## IMPORTANT#############################################

\end{document}